% ch3_e §3.9–3.11 (书页 108–118 / p122–p132)
%--D-TAIL--
%JOIN%Green 函数 (3.71) 沿 $\mathbf{r}$ 和 $\mathbf{r}''$ 方向的球谐函数展开。这些系数依赖于 $\kappa$ 和 $\mathbf{k}$——但除此之外只依赖于晶格结构。实际公式非常复杂，因为其中要涉及各种宗量的 Clebsch–Gordan 系数以及球 Bessel 函数 $j_l$ 和 $n_l$。要计算晶格和，Ewald 方法（§2.3）也很有用（从分母的形式便可以猜到这一点）。不过，一旦针对给定结构编成了数表，系数 $A_{lm;l'm'}$ 就不必再重新计算。

原子势是通过径向函数 $\mathcal{R}_l$ 在球面上的对数导数 $L_l$ 进来的——正如 (3.66) 中那样。之所以如此，是因为尝试函数 (3.73) 在原子球内部满足 Schrödinger 方程，于是我们可以在 (3.70) 与 (3.72) 的被积函数中写出
\begin{equation*}
v_a(\mathbf{r})\,\psi_{\mathbf{k}}(\mathbf{r})=\left(\nabla^2+\mathcal{E}\right)\psi_{\mathbf{k}}(\mathbf{r})\tag{3.76}
\end{equation*}
随后使用 Green 定理，便得到一个面积分，其中出现 Green 函数的导数（它们在 (3.75) 中表现为 $n_l$ 与 $j_l$）以及径向函数的导数。在这个方法中，假定每个原子的 muffin-tin 势是球对称的，这一点至关重要。

由以上讨论可以清楚地看到，Green 函数方法（也称为 Korringa、Kohn 与 Rostoker 方法，即 KKR 方法）与 A.P.W. 方法非常相似。在 A.P.W. 方法中，我们按需要把展开做到任意多的球谐函数项，然后要求解一个久期行列式，以求出来自不同倒格矢的贡献。在 KKR 方法中，我们对格矢求和（实际上，经过对 $G(\kappa,\mathbf{k};\mathbf{r}-\mathbf{r}'')$ 作 Fourier 变换，是对倒格矢求和），但得到的久期行列式是按不同球谐函数的贡献写出的。两种方法都需要高速计算机，而且两者都工作得相当好。KKR 方法的优点是：结构性质与原子性质在行列式的各项中是分离开的，因此可以分别计算、快速合并。A.P.W. 的优点则是：它对原子球以外的波函数给出更好的图像。

\section{模型赝势}

我们在 §3.6 中已经看到，每个原子势对电子的作用，看起来不过是一个弱赝势（pseudo-potential）的作用。这个重要原理，如何体现在基于 muffin-tin 点阵的能带结构技术之中呢？

原子势并不显式地出现在 A.P.W. 和 KKR 公式 (3.65) 与 (3.75) 之中——它们只依赖于 $\mathcal{R}_l$ 在原子球面上的梯度。这个量同样出现在球形势散射的分波理论（partial wave theory）中（见 §5.5）。我们把每个径向解 $\mathcal{R}_l(r,\mathcal{E})$ 与同能量 $\mathcal{E}=\kappa^2$、同角动量 $l$ 的自由电子波匹配起来，从而定义一个\emph{相移}（phase shift）$\eta_l(\mathcal{E})$，它通过下列公式与 $\mathcal{R}_l$ 的对数导数相联系：
\begin{equation*}
L_l\equiv\frac{\mathcal{R}'_l(R_s,\mathcal{E})}{\mathcal{R}_l(R_s,\mathcal{E})}=\left.\frac{j'_l(\kappa r)-\tan\eta_l(\mathcal{E})\cdot n'_l(\kappa r)}{j_l(\kappa r)-\tan\eta_l(\mathcal{E})\cdot n_l(\kappa r)}\right|_{r=R_s}.\tag{3.77}
\end{equation*}
球 Bessel 函数 $j_l$ 与 $n_l$（连同其导数 $j'_l$ 与 $n'_l$）当然正是径向 Schrödinger 方程 (3.60) 在 muffin-tin 势阱之外的自由空间解。

知道了所有分波在所有能量下的相移，我们便得到了关于原子势所需的全部信息；例如在 (3.75) 中，久期行列式的对角元就只含 $\kappa\cot\eta_l(\mathcal{E})$。这十分自然；相移还给出每个原子的散射截面，从而给出晶格对该能量电子的衍射图样。我们可以预期每个原子的行为，就好像它有一个 $t$ 矩阵即\emph{散射振幅}（scattering amplitude），对沿 $\theta$ 角的散射，取标准的分波形式
\begin{equation*}
t_a(\theta)=-\left(\frac{4\pi N}{\kappa}\right)\sum_l(2l+1)\sin\eta_l\exp(i\eta_l)\,P_l(\cos\theta).\tag{3.78}
\end{equation*}
一个很深的势当然可能产生很大的相移——但在 (3.77) 与 (3.78) 中，我们显然可以忽略每个 $\eta_l$ 中 $\pi$ 的所有整数倍。起作用的只是其余部分——而这个余项可以令散射变得非常弱。在分波理论中，当势阱加深时，截面 $|t_a(\theta)|^2$ 并不会急剧增大；随着束缚态的引入它只是起伏波动，很少大于几个电子波长。换言之，Born 近似——其中从 $\mathbf{k}$ 散射到 $\mathbf{k}'$ 的截面正比于 $|\langle\mathbf{k}'|v_a|\mathbf{k}\rangle|^2$——当 $v_a(\mathbf{r})$ 是一个非常深的势时并不成立。这时我们必须用分波公式 (3.78) 取代势的矩阵元。

这正是 §3.2 的简单 N.F.E. 纲要不收敛的原因。我们是在用微扰论去描述晶体中每个原子的散射。但是，为描写波函数在原子内部快速振荡所需要的高阶 Fourier 系数，与芯区之外 Bloch 态的总体衍射图样基本无关。$\mathcal{R}_l$ 的每个节点都给 $\eta_l$ 增加 $\pi$，但对 (3.65)、(3.75) 或 (3.78) 几乎没有影响。因此，O.P.W. 赝势 (3.50) 正是一个解析手段：通过减去正交化的芯态来消去这些内部节点。剩下的弱赝势 $\Gamma$ 只是一个算符，其矩阵元在原子芯之外给出的波函数，与原来的原子势 $v_a$ 所给出的基本相同。

%FIG{60}: 模型势 $w$ 连同原子内部的波函数 $\phi$，重现了真实势 $v$ 连同真实波函数 $\psi$ 的效果。

对赝势概念的这种诠释，为整个能带结构问题提示了一条替代路线：\emph{把每个原子势 $v_a$ 换成一个弱势 $w_a$，它对传导电子具有相同的散射振幅。}直觉上很清楚，原来晶体中的 $\mathcal{E}(\mathbf{k})$ 将与这个假想材料中的能带结构完全相同，而在后者中 N.F.E. 形式体系如今应当是收敛的。我们所要做的，只是把 (3.14) 中的 $\mathcal{V}_{\mathbf{g}'-\mathbf{g}}$，或
(3.64) 中的 $\Gamma_{\mathbf{gg}'}$，换成我们的\emph{模型势}（model potential）$w_a$ 的相应 Fourier 分量。

这一方案会遇到非局域性、能量依赖性和任意性等困难——这些困难在 (3.56) 那类解析赝势中就已被发现。对给定的能量 $\mathcal{E}$ 和角动量 $l$，存在无穷多个局域势 $w_l(\mathbf{r},\mathcal{E})$，在其中求解径向 Schrödinger 方程 (3.60)，都能在原子芯之外精确重现真实的径向函数 $\mathcal{R}_l$。但同一个模型势通常不能同时适用于不同的 $l$ 值或不同的能量，因此我们必须引入一些算符，把波函数中各个角动量分量分离出来，正如 (3.55) 中那样。实践中，\emph{模型赝势}（model pseudo-potential）中每个 $w_l(\mathbf{r},\mathcal{E})$ 的函数形式，与其说是出于深刻的解析理由，不如说主要是为了计算上的方便。

%FIG{61}: 用于 KKR 矩阵元的模型赝势。

在 muffin-tin 点阵中，一个显然的选择是在原子球面上放一个 $\delta$ 函数奇点，对每个 $l$ 值取适当的强度，使之重现散射相移 $\eta_l(\mathcal{E})$（图 61）。利用平面波和球 Bessel 函数的标准解析性质，我们得到如下形式的“赝势”矩阵元：
\begin{equation*}
\Gamma^{\mathrm{KKR}}_{\mathbf{gg}'}=-\frac{4\pi N}{\kappa}\sum_l(2l+1)\tan\eta'_l\,\frac{j_l(|\mathbf{k}-\mathbf{g}|R_s)\,j_l(|\mathbf{k}-\mathbf{g}'|R_s)}{\left\{j_l(\kappa R_s)\right\}^2}\,P_l(\cos\theta_{\mathbf{gg}'}),\tag{3.79}
\end{equation*}
其中
\begin{equation*}
\cot\eta'_l\equiv\cot\eta_l-n_l(\kappa R_s)\big/j_l(\kappa R_s).\tag{3.80}
\end{equation*}
这个公式与 (3.78)——单个孤立原子的散射振幅——十分相似；在区边界附近，当 $|\mathbf{k}-\mathbf{g}|\sim|\mathbf{k}-\mathbf{g}'|\sim\kappa$ 时，在小相移极限下它便趋于后者。

说来也怪，带这些矩阵元的一组 N.F.E. 方程，其实可以从 KKR 方程的久期行列式出发，经过一连串矩阵变换，从角动量表象 (3.64) 变换到倒格子表象而得到。这个公式与 A.P.W. 矩阵元 (3.65) 之间也有代数上的联系——而后者本身是不能从一组简单的局域势导出的。

实践中，对一个给定晶体计算 muffin-tin 势的分布是相当复杂的过程，屏蔽、关联和交换效应在其中起着重要作用（见第 5 章）。非常方便的做法是：在最初阶段——可以说，在把原子聚拢来构成晶格之前——就把每个原子芯内的深势替换为一个模型赝势 $w$。如果让这个赝势在芯外的行为如同价电子所看到的 Coulomb 函数 $-Ze^2/r$，我们甚至可以计算自由原子的光谱能级，并调整 $w$ 的参数去拟合观测值。这一步骤绕开了自洽原子势等许多理论计算，它就是图 62 中勾画的 Heine–Abarenkov 赝势和 Shaw 赝势的基础。在 $w^{\mathrm{HA}}$ 中，‘芯’的内部（半径可任意选取）被做成平坦的，但可以升高或降低，以便对每个角动量都给出正确的电子本征态；在 $w^{\mathrm{S}}$ 中，则为每个 $l$ 值选取一个半径 $R_l$，使得势在‘芯的边缘’处没有急剧的不连续。然而，困难在于要在 Bloch 态的能量 $\mathcal{E}$ 处选取这些参数，而这个能量未必落在光谱能级的集合之内。

%FIG{62}: 模型赝势：(a) Heine–Abarenkov 模型；(b) Shaw 模型。

\section{共振能带}

在紧束缚法（§3.4）的早期应用中，人们注意到 $d$ 轨道如此紧缩在原子芯内部，以至于只通过与近邻的重叠产生很窄的带。在\emph{过渡元素}（transition elements）中，这些内层 $d$ 态并未全部填满，而且离 $s$ 或 $p$ ‘价’态非常近——后者本身组合成一条普通的导带。狭窄的‘d 带’的态密度足以容纳每原子至多 10 个电子，它位于‘s-p 带’之内，并在两者相交处与之\emph{杂化}（hybridize）（图 63）。这种情形当然非常重要，因为伴随着磁性现象。

%FIG{63}: (a) $d$ 带穿过 $s$ 带；(b) $s$–$d$ 杂化。

建立一种内插方案，或者说带可调参数的\emph{模型哈密顿量}（model Hamiltonian），并不困难：其中 $d$ 带用一个 L.C.A.O. 表达式（如 (3.31)）配上适当的赝势分量的矩阵来表示。杂化的程度则由进一步添加的、连接这两个子矩阵的任意矩阵元来确定。

但是，一个像样的定量分析必须顾及 §3.4 中强调的那一点（见图 54）：在我们构成晶体的过程中，原子轨道本身会因势的重叠而被破坏。然而这种破坏并不总是彻底的；在原来 $d$ 轨道的能量附近，仍然可能观察到\emph{虚束缚态}（virtual bound states）。在晶格动力学理论（§2.12）中，我们提到过\emph{共振}（resonance）——一种强烈集中于杂质附近、但在整个晶体中都具有有限振幅的激发。

类似地，电子波函数在共振能量处可以在 muffin-tin 势阱内具有很大的振幅，却并不严格地局域于其中。考察径向 Schrödinger 方程 (3.60) 可知，当 $l\neq0$ 时这种情况很容易发生：‘离心势’ $l(l+1)/r^2$ 提供了一道势垒，电子可以从它的虚态穿过这道势垒隧穿而出（图 64）。

%FIG{64}: 虚态波函数穿透离心势垒的隧穿。

这类共振的存在——不仅来自 $d$ 态，有时甚至来自自由原子的 $p$ 态——是 §3.6 解析赝势形式体系的隐患之一。例如，我们是否应当把这样的态当作‘芯’的一部分来进行正交化，就不清楚。要理解这一局面，必须再一次转向一般散射理论——共振在那里是人们熟悉的现象。在共振能量 $\mathcal{E}_l$ 附近，第 $l$ 个分波的相移按如下公式经过 $\pi/2$：
\begin{equation*}
\tan\eta_l(\mathcal{E})\approx\frac{W_l}{\mathcal{E}-\mathcal{E}_l}.\tag{3.81}
\end{equation*}
把这个公式塞进比如 A.P.W. 或 KKR 程序里，我们只管摇动手柄，就能得到一个能带结构。其结果是在 $\mathcal{E}(\mathbf{k})$ 中产生以 $\mathcal{E}_l$ 为中心的复杂结构，正如紧束缚图像中那样（图 63）。但是，这条\emph{共振能带}（resonance band）在导带中的位置和宽度，如今由 $\mathcal{E}_l$ 的数值——相对于 muffin-tin 零点计量——以及共振的宽度 $W_l$ 决定，而这两者都无法从 L.C.A.O. 模型准确地推出来。

现在若把 (3.81) 代入 KKR 赝势 (3.79)，我们会发现 $\Gamma_{\mathbf{gg}'}$ 在 $\mathcal{E}_l$ 附近出现奇点，把 N.F.E. 方程的收敛性完全破坏掉。这并不是该方法的人为假象。赝势矩阵元强烈的能量依赖性——哪怕在能量上离共振能级相当远——是许多固体的电子能带结构物理的一个本质特征，而在计算中它并不总是得到充分的考虑。

\section{晶体对称性与自旋–轨道相互作用}

在电子能带结构的计算中，晶体结构的对称性十分重要。要正经地讨论它需要群论——这个题目太大，我们在这里无法展开。基本的定理，笼统地说，就是\emph{布里渊区中的能量函数具有晶体的完整点群}。任何使晶体保持不变的操作——比如绕一根轴转动晶体——也同样把函数 $\mathcal{E}(\mathbf{k})$ 变换到它自身。

这可以用来对布里渊区中各点对应的能级和波函数进行分类和化简。为看清楚这是如何发生的，考虑一个具有正方形区的正方格子。取某个像 $\mathbf{k}$ 那样的一般点。假定我们试图把波函数作 A.P.W. 或 O.P.W. 展开，形式为
\begin{equation*}
\psi_{\mathbf{k}}=\sum_{\mathbf{g}}\alpha_{\mathbf{k}-\mathbf{g}}\,\phi_{\mathbf{k}-\mathbf{g}}\tag{3.82}
\end{equation*}
正如 (3.43) 或 (3.63) 那样。我们或许会决定，求和中只需包括图 65($a$) 所示那一组倒格矢。但这样一来，所有矢量 $(\mathbf{k}-\mathbf{g})$ 的大小和方向各不相同，因此所有系数 $\alpha_{\mathbf{k}-\mathbf{g}}$ 也各不相同，彼此之间没有任何明显的关系。我们无法避免去求解一个复杂程度最高的行列式方程。

但假定我们要找的是区对称轴上像 $\Delta$ 这样的某一点处的解。显然，(3.82) 中将会有许多系数相同。
例如，$(\mathbf{g}_5$ 与 $\mathbf{g}_8)$、$(\mathbf{g}_2$ 与 $\mathbf{g}_4)$、$(\mathbf{g}_6$ 与 $\mathbf{g}_7)$ 这些对相对于 $\Delta$ 是对称放置的，因此在展开式中具有相同的系数。于是不必用 8 个各自独立的系数，我们可以假定只有 5 个。这样，行列式的次数较低，求解起来也就更容易。

在对角线上像 $\Sigma$ 这样的点处，情形也一样。在 $X$ 和 $M$ 处还会有进一步的化简，因为这两个点位于倒格矢的中点上，从而可以利用系数之间另外一些等价关系。实际上，区中的点对称性越高，待解方程的次数就越低。\emph{在一定计算量之下，我们在高对称点上能得到比区中任意点处更精确的 $\mathcal{E}(\mathbf{k})$ 值。}

群论告诉我们的还不止这些。它会为我们生成那些具有相同系数的函数——或者，在更复杂的
%FIG{65}: ($a$) 倒格子空间中的一般点 $\mathbf{k}$；($b$) 正方区中的对称点和对称线。
情形下，是乘上像 $1/\sqrt{3}$ 这样一个数后的相同系数。它还告诉我们，在哪些点上可以预期能级之间出现简并。例如，在点 $M$ 处，设有四个 O.P.W. 或 A.P.W. 在能量上简并（对应于图 65 中的 $\mathbf{k}$、$(\mathbf{k}-\mathbf{g}_5)$、$(\mathbf{k}-\mathbf{g}_1)$ 和 $(\mathbf{k}-\mathbf{g}_2)$）。它们会被与晶格势相联系的微扰所劈裂——但可以证明，一个二重简并总会保留下来。当我们画等能面时，这一信息极有帮助。

有一类问题中群论至关重要。在重原子中我们知道存在很强的\emph{自旋–轨道相互作用}（spin–orbit interaction）——一个形如 $\lambda\mathbf{L}\cdot\mathbf{S}$ 的算符，其中 $\mathbf{S}$ 是电子的自旋算符，$\mathbf{L}$ 是轨道角动量算符。在自由原子中，它的作用是消除某些具有同一空间波函数、但电子自旋相反的状态之间的简并。例如，原子的 $P$ 态分裂成 $P_{\frac{3}{2}}$ 和 $P_{\frac{1}{2}}$ 两个态，分别对应于总角动量 $j=\frac{3}{2}$ 和 $j=\frac{1}{2}$。

在构造能带时我们所作的假定是：自旋无关紧要；我们只用一个空间波函数，并且对每个 $\psi_{\mathbf{k}}$ 假定可以把两个自旋相反的电子放进具有这个波函数的两个态里。自旋–轨道相互作用可能解除这些态之间的简并。然而要注意，\emph{时间反演对称性}（time reversal symmetry）这条基本量子原理，总会在任何 Bloch 态 $\psi_{\mathbf{k}}(\mathbf{r})$ 与其复共轭 $\psi_{\mathbf{k}}^*(\mathbf{r})$ 之间保留 \emph{Kramers 简并}（Kramers degeneracy）——后者描写的是电子的波矢和自旋两者都反转了的态。对“无自旋”粒子，这意味着 $\mathcal{E}(\mathbf{k})=\mathcal{E}(-\mathbf{k})$，与晶格的点群对称性无关。

在区中的一般点上，$\psi_{\mathbf{k}}$ 本身不简并；此时若晶格势具有反演中心，自旋–轨道相互作用并不把自旋相反的态分开。这是因为，在态 $\psi_{\mathbf{k}}(\mathbf{r})$ 中反转自旋，等价于考察 $\psi_{\mathbf{k}}(-\mathbf{r})$，而后者具有相同的能量。

但是区中有一些特殊点，自旋–轨道效应在这些点上可能相当重要。最有趣的情形是区的中心，一个具有立方对称性的点。在紧束缚模型中，我们可以用 $p$ 轨道来构造一条‘p 带’。但原子有三个简并的 $p$ 态，所以我们有三条这样的带——比如说，对应于原子的 $p_x$、$p_y$ 和 $p_z$ 轨道。这三条带在 $\mathbf{k}=0$ 处简并；在该点上它们恰好被立方对称性互相变换。每条带的每个态又是二重自旋简并的，所以总共有六重简并。

当我们引入自旋–轨道相互作用时，发现这个六重简并被劈裂成一个四重能级和一个二重能级。实际上，现在我们应当用原子态 $p_{\frac{3}{2}}$——它是四重简并的——和 $p_{\frac{1}{2}}$——它是二重简并的——来构造我们的能带。有趣的是看看当我们沿某个任意方向离开 $\mathbf{k}=0$ 时会发生什么。如果
存在反演对称性，那么每条能级仍然是二重简并的，如图 66($b$) 所示。我们可以说，$j=\frac{3}{2}$ 的带再次分裂成 $m_j=\pm\frac{3}{2}$ 的带（比如说，沿传播方向计量）和 $m_j=\pm\frac{1}{2}$ 的带。如果没有反演对称性，它们还会像图 66($c$) 那样进一步分裂；通过自旋–轨道相互作用，电子自旋变得对晶体势沿传播方向所含的“螺旋”成分敏感了。

%FIG{66}: 自旋–轨道相互作用对区中心附近 $p$ 型能级的影响：($a$) 6 个能级在 $\mathbf{k}=0$ 处简并；($b$) 在具有反演对称性的晶格中，自旋–轨道相互作用使所有态保持二重简并；($c$) 没有反演对称性时，简并被完全解除。

这些效应在诸如 Ge 或 InSb 这类半导体的价带顶部十分重要。在前一种情形，图 66($b$) 成立；在后一种情形，不存在反演中心，区中心处的简并当我们沿任意方向离开时就解除。劈裂的实际符号和大小，可以从自由原子能级中相应效应的符号和大小推测出来。

要作定量计算，方便的办法是把自旋–轨道相互作用 $\lambda\mathbf{L}\cdot\mathbf{S}$ 当作微扰来处理。在区中心附近，它主要作用在 Bloch 函数的周期部分 $u_{\mathbf{k}}$ 上，后者满足方程 (3.38)。事实上，我们可以利用这个方程与普通 Schrödinger 方程的相似性，把附加项
\begin{equation*}
-\frac{i\hbar^2}{m}\,\mathbf{k}\cdot\nabla=\frac{\hbar}{m}\,\mathbf{k}\cdot\mathbf{p}\tag{3.83}
\end{equation*}
当作状态 $u_{\mathbf{0}}$ 上的一个附加微扰（对小的 $\mathbf{k}$ 值），从而不必完整算出整个能带结构，就能细致地计算自旋–轨道劈裂。这就是所谓的 \emph{$\mathbf{k}\cdot\mathbf{p}$ 方法}（$\mathbf{k}\cdot\mathbf{p}$ method）。
